\chapter{Covariance Estimation and Random Matrices}\label{ch:qm:covariance-estimation-and-random-matrices}

A \omterm{def:qm:covariance-estimation-and-random-matrices:pca}{minimum-variance portfolio} of 200 stocks, built from the sample covariance of two years of daily returns, promises an annual volatility of
5.7\%. Held for the next year, it runs at 9.6\%. Nothing broke, and the market did not change: in the chapter's simulation the returns of both years
come from the same factor covariance, whose true \omterm{def:qm:covariance-estimation-and-random-matrices:pca}{minimum-variance portfolio} has a volatility of 7.4\%. With 200 assets and 500 days, the \omterm{def:qm:covariance-estimation-and-random-matrices:cov}{sample
covariance matrix} has 20\,100 parameters estimated from 100\,000 numbers, and the optimiser finds the directions where its errors make risk look
smallest. This chapter is about what a \omterm{def:qm:covariance-estimation-and-random-matrices:cov}{sample covariance matrix} is in high dimension: the law its eigenvalues follow when there is nothing but noise,
the few eigenvalues that carry information, and the \omterm{def:qm:estimation:estimator}{estimators} that clean the rest (shrinkage toward a target, eigenvalue by eigenvalue, clipping,
and factor structure), judged by the one test that matters for a portfolio: realised risk against promised risk.

\section{The sample covariance matrix in high dimension}

\begin{definition}[Covariance matrix, sample covariance matrix]\label{def:qm:covariance-estimation-and-random-matrices:cov}
The \emph{covariance matrix}\index{covariance matrix} of a random vector $r \in \R^N$ is $\Sigma = \E[(r - \mu)(r - \mu)^\top]$. From $T$ observations $r_1, \dots, r_T$
with sample mean $\bar r$, the \emph{sample covariance matrix}\index{sample covariance matrix} is $S = \frac1T\sum_t(r_t - \bar r)(r_t - \bar r)^\top$ (or with divisor
$T - 1$).
\end{definition}

\begin{definition}[Principal component analysis, minimum-variance portfolio]\label{def:qm:covariance-estimation-and-random-matrices:pca}
\emph{Principal component analysis}\index{principal component analysis} diagonalises a covariance (or correlation) matrix, $S = U\Lambda U^\top$: the eigenvectors are
uncorrelated portfolios and the eigenvalues their variances, in decreasing order. The \emph{minimum-variance portfolio}\index{minimum-variance portfolio} for a
covariance $\Sigma$ is $w = \Sigma^{-1}\mathbf 1/(\mathbf 1^\top\Sigma^{-1}\mathbf 1)$, the fully invested portfolio of smallest variance, $1/(\mathbf 1^\top\Sigma^{-1}\mathbf 1)$.
\end{definition}

The \omterm{def:qm:covariance-estimation-and-random-matrices:pca}{minimum-variance portfolio} needs no expected returns, which makes it the cleanest test of a covariance estimate; it also inverts the matrix, which makes
it the harshest. $\Sigma^{-1}$ weights each eigenvector by the reciprocal of its eigenvalue, so the smallest sample eigenvalues, the ones most depressed by
noise, get the most weight. With $q = N/T$ held fixed as both grow, and Gaussian returns, the portfolio built from $S$ has an in-sample variance about $(1 - q)$
times the true minimum and a true variance about $1/(1 - q)$ times it: the realised-to-predicted volatility ratio is about $1/(1 - q)$.

The chapter's truth has 200 stocks, a market factor and five sectors of forty stocks, with specific volatilities between 17\% and 44\% a year. Over fifty
simulated two-year histories ($q = 0.4$), the sample \omterm{def:qm:covariance-estimation-and-random-matrices:pca}{minimum-variance portfolio} predicts 5.71\% and delivers 9.58\% against the truth (9.61\% measured over the next 250
days): a ratio of 1.68 against the formula's $1/(1 - 0.4) = 1.67$. Varying the history from eight years to under one ($q$ from 0.1 to 0.9), the simulated ratio tracks
the formula throughout, reaching 10.1 at $q = 0.9$ (\cref{fig:qm:covariance-estimation-and-random-matrices:bias}).

\begin{omfigure}
\begin{tikzpicture}
\begin{semilogyaxis}[width=10cm,height=5cm,xlabel={$q = N/T$ (200 stocks)},ylabel={true / predicted volatility},
  tick label style={font=\small},label style={font=\small},xmin=0,xmax=0.95,ymin=1,ymax=12,grid=major,grid style={gray!25},
  ytick={1,1.5,2,3,5,10},yticklabels={1,1.5,2,3,5,10},
  legend style={font=\footnotesize,at={(0.03,0.97)},anchor=north west,fill=white,draw=none},table/col sep=comma]
\addplot[only marks,mark=*,mark size=2pt,omDef] table[x=q,y=sim]{figdata/methods/22-covariance-estimation-and-random-matrices/bias.csv};
\addplot[omThm,thick,domain=0:0.93,samples=60] {1/(1-x)};
\legend{simulated (20 histories each),$1/(1 - q)$}
\end{semilogyaxis}
\end{tikzpicture}

\omcaption{Ratio of the true to the predicted volatility of the sample \omterm{def:qm:covariance-estimation-and-random-matrices:pca}{minimum-variance portfolio} of 200 stocks, against $q = N/T$ (histories from 2\,000 to 222 days),
and the high-dimensional formula $1/(1 - q)$. Data: the chapter's tutorial, seeded.}\label{fig:qm:covariance-estimation-and-random-matrices:bias}
\end{omfigure}

\section{The eigenvalue law of pure noise}

\begin{theorem}[Marchenko--Pastur law]\label{thm:qm:covariance-estimation-and-random-matrices:mp}
If the $T \times N$ matrix of returns has iid entries of mean zero and variance $\sigma^2$ and $N/T \to q \in (0, 1)$, the empirical distribution of the eigenvalues of
the \omterm{def:qm:covariance-estimation-and-random-matrices:cov}{sample covariance matrix} converges to the \emph{Marchenko--Pastur law}\index{Marchenko--Pastur law}, with density
\[
\rho(\lambda) = \frac{\sqrt{(\lambda_+ - \lambda)(\lambda - \lambda_-)}}{2\pi q\sigma^2\lambda}, \qquad \lambda_\pm = \sigma^2(1 \pm \sqrt q)^2,
\]
on $[\lambda_-, \lambda_+]$ (Marchenko and Pastur, 1967).
\end{theorem}

The true eigenvalues are all $\sigma^2$; the sample spreads them from $0.14\sigma^2$ to $2.66\sigma^2$ at $q = 0.4$. The smallest are too small by a factor of seven, and they are
the directions a minimum-variance optimiser loves. On 500 days of 200 independent standard normal series, the sample correlation eigenvalues range from 0.144 to 2.565,
inside the theoretical edges 0.135 and 2.665 (\cref{fig:qm:covariance-estimation-and-random-matrices:mp}); Laloux, Cizeau, Bouchaud and Potters (1999) found most of the
eigenvalues of real equity correlation matrices inside the same law.

\begin{omfigure}
\begin{tikzpicture}
\begin{axis}[width=10cm,height=5cm,xlabel={eigenvalue of the sample correlation matrix},ylabel={density},
  tick label style={font=\small},label style={font=\small},xmin=0,xmax=3,ymin=0,ymax=1.2,grid=major,grid style={gray!25},
  legend style={font=\footnotesize,at={(0.97,0.97)},anchor=north east,fill=white,draw=none},table/col sep=comma]
\addplot[omDef,thick,const plot mark mid,fill=omDef!15] table[x=x,y=noise]{figdata/methods/22-covariance-estimation-and-random-matrices/mp_hist.csv};
\addplot[black,very thick] table[x=x,y=density]{figdata/methods/22-covariance-estimation-and-random-matrices/mp_curve.csv};
\addplot[omThm,thick,const plot mark mid] table[x=x,y=factor]{figdata/methods/22-covariance-estimation-and-random-matrices/mp_hist.csv};
\legend{pure noise,Marchenko--Pastur ($q = 0.4$),factor model (bulk)}
\end{axis}
\end{tikzpicture}

\omcaption{Eigenvalues of sample correlation matrices of 200 series over 500 days: independent noise (shaded) against the Marchenko--Pastur density, and the \omterm{def:qm:covariance-estimation-and-random-matrices:factor}{factor model}'s
eigenvalues below 3 (its market eigenvalue, 51.4, and three of its sector eigenvalues lie off the chart). Data: the chapter's tutorial, seeded.}\label{fig:qm:covariance-estimation-and-random-matrices:mp}
\end{omfigure}

\begin{definition}[Spiked covariance model]\label{def:qm:covariance-estimation-and-random-matrices:spiked}
In the \emph{spiked covariance model}\index{spiked covariance model} (Johnstone, 2001) the population covariance is the identity except for a few eigenvalues $\ell_1, \dots, \ell_k
> 1$, the spikes, standing for factors above a noise floor.
\end{definition}

A spike is visible only if it is large enough: when $\ell \le 1 + \sqrt q$ the top sample eigenvalue sticks to the edge $(1 + \sqrt q)^2$; above that threshold it separates and
converges to $\ell(1 + q/(\ell - 1))$, larger than $\ell$ (Baik, Ben Arous and Péché, 2005, for complex data; Baik and Silverstein, 2006, for real). At $q = 0.4$ the threshold is
1.63: a factor explaining less than 1.63 times an average stock's variance is invisible in two years of daily data, however real
(\cref{fig:qm:covariance-estimation-and-random-matrices:spike}). In the \omterm{def:qm:covariance-estimation-and-random-matrices:factor}{factor model}, five sample eigenvalues exceed the edge: the market's 51.4 (population 52.9) and four sector
eigenvalues from 2.96 to 3.68 (population 2.83 to 2.99), each pushed up as the formula says.

\begin{omfigure}
\begin{tikzpicture}
\begin{axis}[width=10cm,height=5cm,xlabel={population spike $\ell$},ylabel={top sample eigenvalue},
  tick label style={font=\small},label style={font=\small},xmin=1,xmax=5.2,ymin=2,ymax=6,grid=major,grid style={gray!25},
  legend style={font=\footnotesize,at={(0.03,0.97)},anchor=north west,fill=white,draw=none},table/col sep=comma]
\addplot[only marks,mark=*,mark size=2pt,omDef] table[x=ell,y=sim]{figdata/methods/22-covariance-estimation-and-random-matrices/spike.csv};
\addplot[omThm,thick,domain=1.633:5.2,samples=60] {x*(1+0.4/(x-1))};
\addplot[omThm,thick,domain=1:1.633,samples=2] {2.665};
\addplot[black,dashed,thin,domain=1:5.2,samples=2] {x};
\legend{{simulated ($N = 200$, $T = 500$)},theory,$y = \ell$}
\end{axis}
\end{tikzpicture}

\omcaption{The spike transition: the largest sample eigenvalue against a single population spike $\ell$ (all other eigenvalues one), with the theory: the bulk edge $(1 + \sqrt q)^2 =
2.66$ below the threshold $1 + \sqrt q = 1.63$, and $\ell(1 + q/(\ell - 1))$ above it. Data: the chapter's tutorial, seeded.}\label{fig:qm:covariance-estimation-and-random-matrices:spike}
\end{omfigure}

\section{Linear and nonlinear shrinkage}

\begin{definition}[Linear shrinkage]\label{def:qm:covariance-estimation-and-random-matrices:linear}
\emph{Linear shrinkage}\index{linear shrinkage} replaces $S$ by $\delta F + (1 - \delta)S$ for a structured target $F$, such as a multiple of the identity or the constant-correlation
matrix, with an intensity $\delta \in [0, 1]$ chosen to minimise the expected Frobenius distance to $\Sigma$.
\end{definition}

It is chapter~14's James--Stein idea applied to a matrix. Ledoit and Wolf (2004) derived consistent estimates of the optimal intensity, which needs only the data. Toward
$mI$, $m = \operatorname{tr}S/N$, it is $\min(\bar b^2, d^2)/d^2$ with $d^2 = \lVert S - mI\rVert^2/N$ and $\bar b^2 = T^{-2}\sum_t\lVert r_tr_t^\top - S\rVert^2/N$, the dispersion of $S$ around the
target against its own sampling noise. On the chapter's histories the intensity toward the identity averages 0.034 and moves the realised volatility from 9.61\% to only
9.14\%: the identity is a poor target when stocks share a market factor. Toward constant correlation the intensity is 0.178, the prediction becomes honest (7.64\%), but the
realised volatility stays at 9.14\%.

\begin{definition}[Rotation-equivariant estimator, nonlinear shrinkage]\label{def:qm:covariance-estimation-and-random-matrices:nonlinear}
A \emph{rotation-equivariant estimator}\index{rotation-equivariant estimator} keeps the sample eigenvectors and changes only the eigenvalues: $\hat\Sigma = U\operatorname{diag}(d)U^\top$.
\emph{Nonlinear shrinkage}\index{nonlinear shrinkage} chooses each $d_i$ separately, as an estimate of the oracle value $u_i^\top\Sigma u_i$, the true variance of the $i$-th
sample eigenportfolio.
\end{definition}

The oracle eigenvalues are the best a \omterm{def:qm:covariance-estimation-and-random-matrices:nonlinear}{rotation-equivariant estimator} can do (they minimise the Frobenius error given $U$). They are unknown, but random-matrix theory expresses
them through the Stieltjes transform of the limiting eigenvalue density, which can be estimated from the sample eigenvalues themselves; Ledoit and Wolf (2020) gave a closed-form
version with a kernel density estimate and its Hilbert transform. On a test with population eigenvalues spread from 0.5 to 2, it cuts the mean squared error of the eigenvalues
against the oracle from 0.49 to 0.0034. On the portfolio test it predicts 7.36\% and realises 8.25\%.

\section{Eigenvalue clipping and factor structure}

\begin{definition}[Eigenvalue clipping]\label{def:qm:covariance-estimation-and-random-matrices:clip}
\emph{Eigenvalue clipping}\index{eigenvalue clipping} (Laloux et al., 1999) keeps the eigenvalues of the sample correlation matrix above the Marchenko--Pastur edge and replaces all
the others by their average, which preserves the trace, then restores the sample variances.
\end{definition}

\begin{definition}[Factor model]\label{def:qm:covariance-estimation-and-random-matrices:factor}
A \emph{factor model}\index{factor model} of covariance is $\Sigma = BF B^\top + D$, with $k \ll N$ factor exposures $B$, a factor covariance $F$ and a diagonal specific variance $D$; statistical
factor models take $B$ from the leading principal components of $S$.
\end{definition}

Clipping predicts 7.27\% and realises 8.05\%; a six-factor principal-component model with a diagonal remainder predicts 6.65\% and realises 7.74\%, closest to the true optimum of 7.44\%
(\cref{fig:qm:covariance-estimation-and-random-matrices:estimators}). The ratios of realised to predicted volatility are 1.68 (sample), 1.51 (linear to the identity), 1.19 (linear to
constant correlation), 1.12 (nonlinear), 1.11 (clipping) and 1.17 (\omterm{def:qm:covariance-estimation-and-random-matrices:factor}{factor model}). The \omterm{def:qm:covariance-estimation-and-random-matrices:factor}{factor model} wins here because the truth is a \omterm{def:qm:covariance-estimation-and-random-matrices:factor}{factor model} with six factors, which the principal
components recover; with a less tidy truth the nonlinear \omterm{def:qm:estimation:estimator}{estimator}, which assumes nothing about structure, is the safer default. Fan, Liao and Mincheva (2013) combine the two: factors
for the spikes, thresholding for the remainder.

\begin{omfigure}
\begin{tikzpicture}
\begin{axis}[width=12.5cm,height=5cm,ybar,bar width=9pt,
  xtick={0,1,2,3,4,5},xticklabels={sample,{LW\\identity},{LW const.\\correlation},nonlinear,clipping,{6-factor\\PCA}},
  xticklabel style={font=\footnotesize,align=center},
  ylabel={annual volatility (\%)},tick label style={font=\small},label style={font=\small},
  xmin=-0.6,xmax=5.6,ymin=0,ymax=11,grid=major,grid style={gray!25},
  legend columns=3,legend style={font=\footnotesize,at={(0.5,-0.3)},anchor=north,draw=none,fill=none,
  /tikz/every even column/.append style={column sep=8pt}},table/col sep=comma]
\addplot[fill=omDef!55,draw=omDef] table[x=k,y=pred]{figdata/methods/22-covariance-estimation-and-random-matrices/estimators.csv};
\addplot[fill=omThm!50,draw=omThm] table[x=k,y=real]{figdata/methods/22-covariance-estimation-and-random-matrices/estimators.csv};
\addplot[black,dashed,thick,sharp plot,forget plot] coordinates {(-0.6,7.44) (5.6,7.44)};
\addlegendimage{line legend,black,dashed,thick}
\legend{predicted (in sample),realised (next year),true optimum}
\end{axis}
\end{tikzpicture}

\omcaption{\omterm{def:qm:covariance-estimation-and-random-matrices:pca}{Minimum-variance portfolios} of 200 stocks from six covariance estimates on two years of daily returns: the volatility each predicts and the volatility each delivers over
the next year, averaged over fifty simulated histories; the dashed line is the true minimum, 7.44\%. Data: the chapter's tutorial, seeded.}\label{fig:qm:covariance-estimation-and-random-matrices:estimators}
\end{omfigure}

\section{Tutorial: the portfolio that promised 6\%}\label{tut:qm:covariance-estimation-and-random-matrices:promise}

\begin{tutorial}
\textbf{Goal.} Build \omterm{def:qm:covariance-estimation-and-random-matrices:pca}{minimum-variance portfolios} from six covariance estimates and compare what they promise with what they deliver.
\textbf{End state:} \cref{fig:qm:covariance-estimation-and-random-matrices:bias,fig:qm:covariance-estimation-and-random-matrices:estimators}; the sample portfolio's ratio of 1.68 and
the better \omterm{def:qm:estimation:estimator}{estimators}' 1.11 to 1.19.
\begin{tutsteps}
\item \textbf{Analytical \omterm{def:qm:covariance-estimation-and-random-matrices:nonlinear}{nonlinear shrinkage}}: kernel density of the eigenvalues, its Hilbert transform, and the shrunk eigenvalues.
\omcode{code/firm/covest/firm_covest.py}{86}{109}{Analytical nonlinear shrinkage of the sample eigenvalues.}\label{lst:qm:covariance-estimation-and-random-matrices:nl}
\item \textbf{Clipping} at the Marchenko--Pastur edge.
\omcode{code/firm/covest/firm_covest.py}{112}{134}{Eigenvalue clipping of the sample correlation matrix.}\label{lst:qm:covariance-estimation-and-random-matrices:clip}
\item \textbf{Run} \texttt{evaluate()}, \texttt{bias\_vs\_q()}, \texttt{spectrum()} and \texttt{spike()} in \texttt{qm\_covest.py}, then \texttt{fig\_covest.py}.
\end{tutsteps}
\textbf{What to change next.} Add a long-only constraint and see how much of the sample portfolio's error it removes (it acts as a shrinkage); replace Gaussian returns by Student-$t$
returns and compare the \omterm{def:qm:estimation:estimator}{estimators} again.
\end{tutorial}

\section{Build: the covariance estimators}\label{bld:qm:covariance-estimation-and-random-matrices:covest}

\begin{build}
\bfield{Purpose} Every covariance the miniature firm's portfolio construction and risk use (chapter~23 onward) comes from this module, and its realised-to-predicted risk ratio is
tracked.
\bfield{Interface} \texttt{sample\_cov(X)}; \texttt{mp\_edges(q)}, \texttt{mp\_density(x, q)}; \texttt{lw\_identity(X)}, \texttt{lw\_constant\_corr(X)}; \texttt{nonlinear\_shrinkage(X)};
\texttt{clip(X)}; \texttt{pca\_factor(X, k)}; \texttt{min\_var\_weights(Sigma)}; \texttt{portfolio\_risk(w, Sigma)}.
\bfield{Rules} Never invert a raw sample covariance when $N/T$ exceeds a few percent; report the intensity or the clipping threshold with every estimate; judge \omterm{def:qm:estimation:estimator}{estimators} by out-of-sample
risk of the portfolios built from them, not by in-sample fit.
\bfield{Acceptance tests} \texttt{code/firm/covest/tests/}: the Marchenko--Pastur density integrates to one and bounds the noise spectrum; the Ledoit--Wolf intensity matches its brute-force
definition and improves the Frobenius error; the constant-correlation target improves on a constant-correlation truth; \omterm{def:qm:covariance-estimation-and-random-matrices:nonlinear}{nonlinear shrinkage} tracks the oracle eigenvalues; clipping and the
\omterm{def:qm:covariance-estimation-and-random-matrices:factor}{factor model} keep the sample variances; minimum-variance weights for a diagonal covariance.
\bfield{Stretch} POET (factors plus thresholding); \omterm{def:qm:covariance-estimation-and-random-matrices:nonlinear}{nonlinear shrinkage} for $N > T$; a rolling realised-to-predicted monitor on live portfolios.
\end{build}

\begin{omsources}
\item V.~A.~Marčenko and L.~A.~Pastur, ``Distribution of eigenvalues for some sets of random matrices'', \emph{Mathematics of the USSR-Sbornik} 1, 1967.
\item L.~Laloux, P.~Cizeau, J.-P.~Bouchaud and M.~Potters, ``Noise dressing of financial correlation matrices'', \emph{Physical Review Letters} 83, 1999.
\item O.~Ledoit and M.~Wolf, ``A well-conditioned estimator for large-dimensional covariance matrices'', \emph{Journal of Multivariate Analysis} 88, 2004; ``Honey, I shrunk the sample
covariance matrix'', \emph{Journal of Portfolio Management} 30, 2004.
\item O.~Ledoit and M.~Wolf, ``Analytical nonlinear shrinkage of large-dimensional covariance matrices'', \emph{Annals of Statistics} 48, 2020.
\item I.~M.~Johnstone, ``On the distribution of the largest eigenvalue in principal components analysis'', \emph{Annals of Statistics} 29, 2001.
\item J.~Baik, G.~Ben Arous and S.~Péché, \emph{Annals of Probability} 33, 2005; J.~Baik and J.~W.~Silverstein, \emph{Journal of Multivariate Analysis} 97, 2006.
\item J.~Fan, Y.~Liao and M.~Mincheva, ``Large covariance estimation by thresholding principal orthogonal complements'', \emph{Journal of the Royal Statistical Society B} 75, 2013.
\item J.~Bun, J.-P.~Bouchaud and M.~Potters, ``Cleaning large correlation matrices: tools from random matrix theory'', \emph{Physics Reports} 666, 2017.
\end{omsources}

\section{Exercises}

\begin{exercise}[$\star$]\label{exo:qm:covariance-estimation-and-random-matrices:1}
What are the Marchenko--Pastur edges for $N = 500$ stocks and $T = 1\,250$ days?
\end{exercise}

\begin{exercise}[$\star$]\label{exo:qm:covariance-estimation-and-random-matrices:2}
A sample \omterm{def:qm:covariance-estimation-and-random-matrices:pca}{minimum-variance portfolio} of 300 stocks built from 600 days predicts 5\%. What volatility should one expect?
\end{exercise}

\begin{exercise}[$\star$]\label{exo:qm:covariance-estimation-and-random-matrices:3}
Show that the \omterm{def:qm:covariance-estimation-and-random-matrices:pca}{minimum-variance portfolio}'s variance is $1/(\mathbf 1^\top\Sigma^{-1}\mathbf 1)$.
\end{exercise}

\begin{exercise}[$\star\star$]\label{exo:qm:covariance-estimation-and-random-matrices:4}
With $q = 0.25$, which population spikes are detectable, and where does a spike of 3 appear in the sample?
\end{exercise}

\begin{exercise}[$\star\star$]\label{exo:qm:covariance-estimation-and-random-matrices:5}
Show that the \omterm{def:qm:covariance-estimation-and-random-matrices:linear}{linear shrinkage} $\delta mI + (1 - \delta)S$ keeps the eigenvectors of $S$ and maps each eigenvalue $\lambda$ to $\delta m + (1 - \delta)\lambda$. Why is this less flexible than
\omterm{def:qm:covariance-estimation-and-random-matrices:nonlinear}{nonlinear shrinkage}?
\end{exercise}

\begin{exercise}[$\star\star$]\label{exo:qm:covariance-estimation-and-random-matrices:6}
How many days of data would make the sample \omterm{def:qm:covariance-estimation-and-random-matrices:pca}{minimum-variance portfolio} of 200 stocks underpredict its volatility by at most 10\%?
\end{exercise}

\begin{exercise}[$\star\star\star$]\label{exo:qm:covariance-estimation-and-random-matrices:7}
\emph{Coding.} Reproduce the ratio of true to predicted volatility of the sample \omterm{def:qm:covariance-estimation-and-random-matrices:pca}{minimum-variance portfolio} for $q$ from 0.1 to 0.9 and compare it with $1/(1 - q)$.
\end{exercise}

\begin{exercise}[$\star\star\star$]\label{exo:qm:covariance-estimation-and-random-matrices:8}
\emph{Find the flaw.} ``We chose our covariance \omterm{def:qm:estimation:estimator}{estimator} because it has the lowest in-sample \omterm{def:qm:covariance-estimation-and-random-matrices:pca}{minimum-variance portfolio} volatility of all the \omterm{def:qm:estimation:estimator}{estimators} we tried.''
\end{exercise}

\section{Problem: The Portfolio That Promised 6\%}

\begin{problem}[{Weekend problem --- risk promised and risk delivered}]\label{pb:qm:covariance-estimation-and-random-matrices:1}
Two hundred stocks follow a known covariance with a market factor and five sectors of forty. A manager estimates the covariance from two years (500 days) of daily returns, builds the
\omterm{def:qm:covariance-estimation-and-random-matrices:pca}{minimum-variance portfolio} and holds it for a year. Fifty histories are simulated.

\medskip\noindent\textbf{Part I --- The sample.}
\begin{enumerate}
\item What does the sample portfolio predict and deliver, and what is the true optimum?
\item What ratio does the formula $1/(1 - q)$ predict, and what does the simulation give?
\item What are the Marchenko--Pastur edges at $q = 0.4$, and where do pure-noise eigenvalues fall?
\item Which eigenvalues of the \omterm{def:qm:covariance-estimation-and-random-matrices:factor}{factor model} stand out, and by how much are they biased?
\item What is the smallest detectable factor at this $q$?
\end{enumerate}

\medskip\noindent\textbf{Part II --- Better \omterm{def:qm:estimation:estimator}{estimators}.}
\begin{enumerate}[resume]
\item What intensities do the two \omterm{def:qm:covariance-estimation-and-random-matrices:linear}{linear shrinkages} choose, and what do they deliver?
\item What do \omterm{def:qm:covariance-estimation-and-random-matrices:nonlinear}{nonlinear shrinkage} and clipping deliver?
\item What does the six-factor model deliver, and why is it best here?
\item What are the realised-to-predicted ratios of all six?
\item Which \omterm{def:qm:estimation:estimator}{estimator}'s prediction is most honest, and which portfolio is best?
\end{enumerate}

\medskip\noindent\textbf{Part III --- Dimensions.}
\begin{enumerate}[resume]
\item What happens to the sample portfolio at $q = 0.9$?
\item How many days would make the ratio 1.1?
\item Why does the optimiser pick the noisiest directions?
\item What does a long-only constraint do to this problem?
\item Why does the identity make a poor shrinkage target for stocks?
\end{enumerate}

\medskip\noindent\textbf{Part IV --- Judgement.}
\begin{enumerate}[resume]
\item Which \omterm{def:qm:estimation:estimator}{estimator} would you put in production, and why?
\item How would you monitor it?
\item What would change with fat-tailed returns or a changing covariance?
\item State the \emph{named result}: the realised-to-predicted risk ratio of the sample portfolio at $q = 0.4$ against $1/(1 - q)$, and the ratio each better \omterm{def:qm:estimation:estimator}{estimator} achieves.
\item In one sentence: why does an optimiser make estimation error worse?
\end{enumerate}
\end{problem}

\section{Interview questions}

\begin{interviewq}[$\star$ \iqroles{researcher, risk}]\label{iq:qm:covariance-estimation-and-random-matrices:1}
Why does a \omterm{def:qm:covariance-estimation-and-random-matrices:pca}{minimum-variance portfolio} built from a sample covariance underpredict its risk?
\end{interviewq}

\begin{interviewq}[$\star\star$ \iqroles{researcher, mle}]\label{iq:qm:covariance-estimation-and-random-matrices:2}
What is the \omterm{thm:qm:covariance-estimation-and-random-matrices:mp}{Marchenko--Pastur law}, and how do you use it on a correlation matrix?
\end{interviewq}

\begin{interviewq}[$\star\star$ \iqroles{researcher}]\label{iq:qm:covariance-estimation-and-random-matrices:3}
Explain Ledoit--Wolf shrinkage. What does the intensity depend on?
\end{interviewq}

\begin{interviewq}[$\star\star$ \iqroles{researcher, risk}]\label{iq:qm:covariance-estimation-and-random-matrices:4}
You have 1\,000 stocks and three years of daily data. How do you estimate the \omterm{def:qm:covariance-estimation-and-random-matrices:cov}{covariance matrix}?
\end{interviewq}

\begin{interviewq}[$\star\star$ \iqroles{mle, researcher}]\label{iq:qm:covariance-estimation-and-random-matrices:5}
How many principal components of a stock correlation matrix carry information?
\end{interviewq}

\begin{interviewq}[$\star\star\star$ \iqroles{researcher}]\label{iq:qm:covariance-estimation-and-random-matrices:6}
What is a \omterm{def:qm:covariance-estimation-and-random-matrices:nonlinear}{rotation-equivariant estimator}, and what is the best one could do in that class?
\end{interviewq}
